% !TeX root = ../../thesis.tex
\markboth{\uppercase{Resume sommaire}}{\uppercase{Resume sommaire}}
\vspace{-0.9cm}
\chapter{Résumé sommaire}
\markright{\uppercase{Resume sommaire}}
\instructionsabstract
\vspace{-0.6cm}
%La régulation traductionnelle reposant sur l'usage de codons synonymes joue un rôle important en biologie moléculaire et notamment dans la compréhension des mécanismes de résistance aux therapies ciblées contre les cancers. Les ARN de transfert (ARNt) modifiés sont des acteurs clés impliqués dans la régulation des niveaux d'expression des protéines en modulant différemment le décodage de certains codons synonymes. Néanmoins, leur contribution à la dynamique de la synthèse protéique reste mal comprise.

%Les modèles basés sur des agents (ABM), tels que les modèles du procédé d'exclusion simple totalement asymétrique (TASEP), ont été utilisés pour quantifier les taux de traduction des transcrits par les ribosomes. Notre travail vise à étendre la modélisation ABM et TASEP afin de prendre en compte les effets des modifications des ARN de transfert. Nous avons mis en œuvre un modèle computationnel stochastique quantifiant les vitesses de synthèse protéique.

%L'algorithme utilise le temps de résidence des ribosomes sur chaque codon étant données les séquences de codons des transcrits, l'abondance relative des transcrits ainsi que des tables des cinétiques d'accommodation des ARN de transfert ARNt (modifiés ou non), de formation des liaisons peptidiques et de translocation. Des caractéristiques importantes du modèle incluent la variation du taux d'élongation causée par les acides aminés chargés dans le tunnel de sortie du ribosome, la transpeptidation de la proline au centre de transfert peptidique et les effets de ralentissement induits par la structure secondaire des transcrits.

%Le modèle permet de comparer les niveaux d'expression protéique relatifs ainsi que les profils de polysomes et de RiboSeq dans différents scénarios, avec un pool de ribosomes et des taux d'initiation paramétrables. Notre intention est d'utiliser notre modèle pour mieux comprendre comment l'usage des codons et les modifications des ARNt interagissent dynamiquement et influencent la synthèse protéique.

%L'objectif ultime de ce travail est de comprendre comment le contrôle traductionnel via l'utilisation des codons et l'ajustement des ARNt influence le protéome, la structure locale des protéines, et potentielement les changements d'activité fonctionnelle spécifique, l'agrégation, la ribophagie, l'autophagie ou la dégradation protéasomale. Les résultats de cette étude pourraient contribuer au développement de thérapies ciblées contre les cancers, notamment le mélanome et les cancers du poumon.

%---------------------------------

La régulation traductionnelle reposant sur l'usage de codons synonymes joue un rôle important en biologie moléculaire et notamment dans la compréhension des mécanismes de résistance aux therapies ciblées contre les cancers. Les ARN de transfert (ARNt) modifiés sont des acteurs essentiels impliqués dans la régulation des niveaux d'expression protéique en optimisant le décodage des codons utilisés de manière différentielle. Néanmoins, leur contribution à la dynamique de la synthèse protéique reste encore mal comprise.

Les modèles multi-agents (Agent-Based Models, ABMs), en particulier le processus d'exclusion totalement asymétrique (Totally Asymmetric Simple Exclusion Process, TASEP), ont été employés pour quantifier les taux de traduction par les ribosomes. En s’appuyant sur ces fondations, notre travail étend les modèles ABM et TASEP afin d’intégrer les effets des modifications des ARNt ainsi que d'autres facteurs influençant la traduction. Nous avons développé un modèle stochastique computationnel permettant de quantifier les taux de synthèse protéique, en utilisant les temps de résidence des ribosomes sur chaque codon, basés sur les séquences des transcrits et leur abondance relative. Le modèle intègre, par codon et pour chaque ribosome individuel, les taux spécifiques d’accommodation des ARNt (modifiés ou non), la formation de la liaison peptidique, et la translocation.

Les caractéristiques clés du modèle incluent la modulation des vitesses d’élongation par les acides aminés chargés dans le tunnel de sortie du ribosome, la transpeptidation de la proline au centre de la peptidyl transférase, ainsi que le ralentissement de la traduction dû aux structures secondaires de l'ARNm. Notre modèle permet la comparaison des niveaux relatifs d’expression protéique, des profils de polysomes, et des données de type Ribo-Seq dans divers scénarios, en tenant compte de la disponibilité du pool ribosomal et des taux d'initiation.

En définitive, ce modèle vise à élucider comment l’utilisation des codons et les modifications des ARNt interagissent dynamiquement pour influencer la synthèse protéique. En démêlant systématiquement les facteurs contributifs, nous cherchons à évaluer leur sensibilité non seulement sur la production protéique, mais également sur la densité ribosomale et les profils de polysomes.

%%%%%%%%%%%%%%%%%%%%%%%%%%%%%%%%%%%%%%%%%%%%%%%%%%
% Keep the following \cleardoublepage at the end of this file,
% otherwise \includeonly includes empty pages.
\cleardoublepage

% vim: tw=70 nocindent expandtab foldmethod=marker foldmarker={{{}{,}{}}}
